\section{The scalar unordered-pair representation is also excluded}
\label{sec:unordered-pair-obstruction}

\paragraph{Scope and check status.}
This is a new elementary proof, recorded on 21 November 2026.
It concerns the representation usually denoted $S^{(n-2,2)}$,
of dimension $n(n-3)/2$, and scalar twists of its restrictions.
No hypothesis on the scalar character field is imposed.
The finite check for $5\le n\le10$ has run exactly: all 127
partitions and all 751 matching scalar phases pass, with a separate
check for every off-spectrum scalar.  Two different exact spectral
constructions agree.  The parent completed an independent whole-proof
review of D-UP-01--03, including the all-degree inequalities, arbitrary-phase
exhaustiveness, projection descent, and the degree-four exception, and
recorded PASS.  A second independent whole-proof review by the G34
agent also passed; its separate permutation-cycle implementation
reconstructed all 127 partitions and 751 matching phases exactly.

\subsection*{D-UP-01. Representation and spectral estimates}

Let $P_n$ be the permutation space on $n$ letters, let
$U_n=\operatorname{Std}_n$ be its deleted summand, and let $P_{2,n}$
be the permutation space on unordered pairs of distinct letters.
The incidence map
\[
        e_i\longmapsto\sum_{j\ne i}e_{\{i,j\}}
\]
embeds $P_n$ into $P_{2,n}$ for $n\ge3$: a relation between its
columns would have coefficients satisfying $a_i+a_j=0$ for all
distinct $i,j$, and hence all $a_i=0$.
Let $W_{2,n}$ be its orthogonal complement.  Thus
\[
  P_{2,n}=P_n\oplus W_{2,n},\qquad
  m=\dim W_{2,n}=\frac{n(n-3)}2.                       \tag{UP.1}
\]
Everything is defined over $\mathbf Q$ and is real.
The identity
$\operatorname{Sym}^2P_n=P_n\oplus P_{2,n}$, together with
$P_n=1\oplus U_n$, also gives
\[
          \operatorname{Sym}^2U_n\simeq P_{2,n}.
                                                               \tag{UP.2}
\]
These identities follow by decomposing diagonal and off-diagonal
monomials and cancelling the common summand in characteristic zero.

Fix $n\ge5$.  For a permutation $\pi$, write $c$ for its number
of cycles and $k=n-c$.  For a root of unity $\lambda$, set
$D_\lambda=\operatorname{rank}(1-\lambda W_{2,n}(\pi))$.
We prove
\begin{align}
 \pi\ne1&\quad\Longrightarrow\quad
 D_1\ge\frac{n-5}{2}k+\frac n4,                       \tag{UP.3}\\
 \lambda\ne1&\quad\Longrightarrow\quad
 D_\lambda\ge\frac{n-5}{2}k+\frac{n-4}{2}.             \tag{UP.4}
\end{align}

\paragraph{Proof of (UP.3).}
Let $l_1,\ldots,l_c$ be the cycle lengths, let $e$ count the even
ones, and put
\[
    Q=\sum_{i<j}\bigl(l_i+l_j-2\gcd(l_i,l_j)\bigr)\ge0.
\]
Two different cycles contribute $\gcd(l_i,l_j)$ fixed unordered-pair
orbits.  One cycle contributes $\lfloor l_i/2\rfloor$, including
its opposite-pair orbit when $l_i$ is even.  Subtracting the $c$
fixed vectors of $P_n$ in (UP.1) gives
\[
                     D_1=\frac{(n-3)k-e+Q}{2}.         \tag{UP.5}
\]
Equivalently it is the exterior-square moving rank of (EX.5) minus
the number of even cycles.
We claim $2k-e+Q\ge n/2$ whenever $\pi\ne1$.
If there are no fixed points, then $k\ge n/2$ and $e\le k$,
which proves the claim.  If there are $f\ge1$ fixed points, then
$Q\ge fk$, $e\le k$, $k\ge1$, and $n\le f+2k$.
Consequently
\[
       2k-e+Q\ge(f+1)k\ge\frac{f+2k}{2}\ge\frac n2.
\]
The middle inequality is $f(2k-1)\ge0$.
Subtracting $(n-5)k/2$ from (UP.5) now proves (UP.3).

\paragraph{Proof of (UP.4).}
Put $t=\lambda^{-1}\ne1$.  On $U_n$, the eigenvalue $1$ has
multiplicity $a=c-1$, the complementary space has dimension $k$,
and each eigenvalue has multiplicity at most $c$.
Let $b$ be the multiplicity of $t$ on that complement.
In $\operatorname{Sym}^2U_n$, the pair of eigenvalue spaces $1,t$
contributes $ab$.  An eigenvector with eigenvalue $t$ cannot
contribute to any other symmetric pair of product $t$.
The remaining relevant space therefore has dimension $k-b$.
On it, two distinct eigenvalue spaces of dimensions $x,y\le c$
contribute at most $c(x+y)/2$, and a self-paired eigenspace of
dimension $x$ contributes $\binom{x+1}{2}\le(c+1)x/2$.
By (UP.1)--(UP.2) one must also subtract the $b$ eigenvectors of
$P_n$ with eigenvalue $t$.  Hence the fixed dimension $F$ of
$\lambda W_{2,n}(\pi)$ is bounded by
\[
 F\le ab-b+\frac{c+1}{2}(k-b)
       =\frac{(c+1)k+(c-5)b}{2}.                       \tag{UP.6}
\]
If $c\ge5$, use $b\le c$.  Substituting in $m-F$ gives
\[
            D_\lambda-\frac{n-5}{2}k\ge\frac{n+c}{2}.
\]
If $1\le c\le4$, use $b\ge0$ and $c-5<0$ instead.  It gives
\[
 D_\lambda-\frac{n-5}{2}k
       \ge\frac{n+c(c-4)}2\ge\frac{n-4}{2},
\]
because $c(c-4)\ge-4$ for these four integers.  This proves (UP.4).

Both right-hand sides are positive in their stated nonidentity cases.
It follows in particular that $W_{2,n}(\mathfrak S_n)$ intersects
the scalar matrices only in the identity for $n\ge5$.
Thus, for a representation $\chi\otimes(W_{2,n}\circ\pi)$,
both $\pi$ and $\chi$ descend to its faithful image.
This conclusion is not asserted at $n=4$.

\subsection*{D-UP-02. Uniform obstruction and the finite degrees}

Consider an irreducible scalar twist
\[
           \tau(j)=\chi(j)W_{2,n}(\pi(j)),\qquad n\ge5,
                                                               \tag{UP.7}
\]
whose character is nonreal.  Suppose a genus-zero branch tuple in
its faithful image has product one, generates that image, and has
six-dimensional cohomological multiplicity.  Then
\[
                    \sum_jD_j=2m+6.                    \tag{UP.8}
\]
The permutation image is transitive; the verification from
irreducibility is given in D-UP-03 below.  Because the projections
descend by D-UP-01, the product-one tuple gives an actual transitive
degree-$n$ permutation cover of the sphere.  Therefore
\[
                    \sum_j\operatorname{ind}(\pi_j)\ge2n-2.
                                                               \tag{UP.9}
\]
Its scalar quotient is nonreal.  There are at least two nonreal
scalar entries, since one nonreal scalar cannot have product one
with any number of real roots of unity.  In particular the number
$s$ of entries with $\lambda_j\ne1$ satisfies $s\ge2$.

Suppose $n\ge11$ and $r\ge12$.  Applying (UP.3)--(UP.4) gives
\begin{align*}
 \sum_jD_j
 &\ge \frac{n-5}{2}\sum_j\operatorname{ind}(\pi_j)
          +\frac n4(r-s)+\frac{n-4}{2}s\\
 &\ge (n-5)(n-1)+3n+\frac{n-8}{4}s\\
 &\ge n^2-3n+5+\frac{n-8}{2}\\
 &=2m+6+\frac{n-10}{2}>2m+6.
\end{align*}
This contradicts (UP.8).  The strict inequality already holds for
$n=11$; no limiting or asymptotic argument is used.

For $5\le n\le10$, the exact minimum of $D_\lambda(\pi)$ over
all genuine scalar-permutation matrices, allowing arbitrary scalar
orders, is $n-3$.  The complete finite certificate is as follows.
For a partition $l_1+\cdots+l_c=n$, form a multiset of eigenvalue
fractions for $P_{2,n}$:
\begin{itemize}
 \item two different cycles of lengths $a,b$ contribute
 $\gcd(a,b)$ copies of every fraction $j/\operatorname{lcm}(a,b)$;
 \item one cycle of length $l$ contributes
 $\lfloor(l-1)/2\rfloor$ copies of every fraction $j/l$;
 \item if $l$ is even, add one copy of every fraction $j/(l/2)$
 for its opposite-pair orbit.
\end{itemize}
Subtract one copy of every fraction $j/l_i$ for each cycle to
remove $P_n$.  All resulting multiplicities are nonnegative and
sum to $m$, by (UP.1).  If $q$ occurs with multiplicity $b_q$,
the scalar $\exp(-2\pi iq)$ gives moving dimension $m-b_q$.
Every scalar not of this form gives moving dimension $m$.
Thus this finite list tests every possible scalar order.
Omit only $\pi=1$, $q=0$, which is the identity branch matrix.

The accompanying program enumerates the partitions recursively by
their nondecreasing parts, uses exact fractions, and compares this
orbit multiset to the independently constructed multiset of all
unordered sums with repetition of standard-representation
eigenvalue fractions, followed by subtraction of $P_n$.
The latter construction is (UP.2).
The complete results are
\[
\begin{array}{c|rrrrrr}
 n&5&6&7&8&9&10\\ \hline
 m&5&9&14&20&27&35\\
 \#\text{partitions}&7&11&15&22&30&42\\
 \min D_\lambda&2&3&4&5&6&7\\
 12\min D_\lambda&24&36&48&60&72&84\\
 2m+6&16&24&34&46&60&76
\end{array}                                                     \tag{UP.10}
\]
Every minimum is attained by a scalar-$1$ transposition.
The strict inequalities in the last two rows exclude $r\ge12$.
The certificate executed 127 partitions and 751 matching scalar
checks; off-spectrum scalars were treated as just explained.
This is an exhaustive finite proof of the small-degree inequalities,
not a numerical sampling of partitions or phases.

\subsection*{D-UP-03. Inducing representations and the exceptional degree four}

Let an exact-imaginary-quadratic-field Weil sixfold representation
be induced from an irreducible representation of the form (UP.7),
now allowing $n\ge4$.  Then its period rank is at most eight.
There is no assumption that the inducing character field or scalar
field equals the original quadratic field.  In particular none of
these blocks supplies a dominant family, split or nonsplit.

First let $n\ge5$.  We verify the transitivity used above.
If the permutation image preserves a subset of size $a$ with
$2\le a\le n-2$, it lies in
$\mathfrak S_a\times\mathfrak S_{n-a}$.
This full subgroup has three orbits on unordered pairs and two
orbits on letters.  Hence (UP.1) gives a one-dimensional invariant
subspace of $W_{2,n}$, a proper subspace since $m\ge5$.
It stays invariant under the smaller permutation image and under
the scalar twist, contrary to irreducibility.
If no such subset exists but the image is intransitive, it fixes
one letter and is transitive on the remaining $n-1$ letters.
Restricting (UP.1) to its containing $\mathfrak S_{n-1}$ gives
\[
 W_{2,n}|_{\mathfrak S_{n-1}}
       \simeq U_{n-1}\oplus W_{2,n-1}.                  \tag{UP.11}
\]
Indeed $P_{2,n}$ restricts as $P_{2,n-1}\oplus P_{n-1}$,
whereas $P_n$ restricts as $1\oplus P_{n-1}$.
Both summands in (UP.11) have positive dimension for $n\ge5$.
They remain invariant after restriction and scalar twist, again
contradicting irreducibility.  Thus the permutation image is
transitive.

Pass to the intermediate connected cover for the faithful image
of the inducing representation.  Its complex multiplicity
variation is the original one by Frobenius reciprocity and
quotient pullback, and the fixed Hermitian polarization determines
the conjugate filtration.  The full quadratic-field period rank
therefore factors through this intermediate-cover moduli, as in
D-BI-01 and D-EX-03.  The scalar and permutation maps descend by
(UP.3)--(UP.4), so no projected branch data have been lost.
If $h$ is the genus of the intermediate base, its Euler equation is
\[
                       6=2m(h-1)+\sum_jD_j.
\]
For $h=1$ the marked-base moduli has dimension at most six.
For $h\ge2$ the right-hand contribution $2m(h-1)$ already
exceeds six because $m\ge5$.
Thus rank nine would require $h=0$ and $r\ge12$.
The inducing character is nonreal: otherwise its induced character
would be real, contrary to the original exact imaginary quadratic
field.  Because $W_{2,n}$ is rational, a scalar image contained in
$\{1,-1\}$ would give a real inducing character.  Its scalar image
is therefore nonreal, and D-UP-02 supplies the contradiction.

When $n=4$, the representation $W_{2,4}$ has dimension two and
need not retain a faithful degree-four permutation projection.
We use no such projection in this case.
For any irreducible inducing representation of dimension two with
cohomological multiplicity six, the Euler equation gives
$D=10-4h$.  Since each genuine branch contributes at least one,
the dimensions of its marked-base moduli are at most seven for
$h=0$, six for $h=1$, and five for $h=2$; larger $h$ is impossible.
Thus the period rank is at most seven.  This treats degree four
without an incorrect projective-faithfulness assertion.

These obstructions use only necessary local spectra, the Euler
equation, and Riemann--Hurwitz.  No feasible aggregate profile is
being promoted to a product-one generating tuple, and no source
dimension inequality is being promoted to dominance.
